\section{Aplikasi dan Studi Kasus}

Bagian ini menyajikan dua studi kasus yang mengaitkan aturan inferensi dan metode bukti dengan sistem berbasis aturan serta verifikasi spesifikasi program \cite{rosen2019,stanford_cs103,mit_6042}.

\subsection*{Studi Kasus 1: Sistem Aturan Sederhana dengan Modus Ponens dan Modus Tollens}

Sistem berbasis aturan (\emph{rule-based}) banyak dipakai dalam sistem pakar, diagnosa, atau rekomendasi. Setiap aturan berbentuk ``jika kondisi A maka kesimpulan B'' ($A \to B$). Diberikan fakta (premis), kita menarik kesimpulan dengan \textbf{modus ponens}: jika aturan $A \to B$ ada dan fakta $A$ benar, maka $B$ benar. Sebaliknya, \textbf{modus tollens} dipakai ketika kita mengetahui $\neg B$: maka $\neg A$ harus benar (karena jika $A$ benar, $B$ akan benar, kontradiksi dengan $\neg B$) \cite{rosen2019,levin_discrete}. Contoh: aturan ``jika hari hujan maka jalan licin''; jika kita tahu ``jalan tidak licin'', maka dengan modus tollens disimpulkan ``tidak hujan''.

\begin{contoh}
\textbf{Mesin inferensi sederhana dengan modus ponens dan modus tollens}

Program berikut menyimpan aturan implikasi dan fakta, lalu menerapkan modus ponens dan modus tollens untuk menambah kesimpulan. Aturan dinyatakan sebagai pasangan (anteseden, konsekuen); fakta adalah himpunan proposisi yang diketahui benar.

\begin{lstlisting}[style=pythonstyle, caption={Sistem aturan: modus ponens dan modus tollens untuk menarik kesimpulan}]
def implikasi(p, q):
    """p -> q (False hanya jika p True dan q False)."""
    return (not p) or q

def inferensi_modus_ponens(aturan, fakta):
    """Dari aturan (p, q) dan fakta p, tambahkan q ke fakta."""
    kesimpulan_baru = set()
    for (p, q) in aturan:
        if p in fakta and q not in fakta:
            kesimpulan_baru.add(q)
    return fakta | kesimpulan_baru

def inferensi_modus_tollens(aturan, fakta):
    """Dari aturan (p, q) dan fakta not q, tambahkan not p."""
    kesimpulan_baru = set()
    for (p, q) in aturan:
        not_q = ("not", q)
        not_p = ("not", p)
        if not_q in fakta and not_p not in fakta:
            kesimpulan_baru.add(not_p)
    return fakta | kesimpulan_baru

# Aturan: hujan -> jalan_licin; jalan_licin -> hati_hati
aturan = [("hujan", "jalan_licin"), ("jalan_licin", "hati_hati")]
fakta = {"hujan"}
fakta = inferensi_modus_ponens(aturan, fakta)  # tambah jalan_licin
fakta = inferensi_modus_ponens(aturan, fakta)  # tambah hati_hati
print("Setelah modus ponens:", fakta)  # {hujan, jalan_licin, hati_hati}

# Modus tollens: fakta (not, jalan_licin) -> simpulkan (not, hujan)
fakta2 = {("not", "jalan_licin")}
# Dengan aturan hujan -> jalan_licin, maka not jalan_licin -> not hujan
print("Contoh modus tollens: tidak jalan licin -> tidak hujan")
\end{lstlisting}
\end{contoh}

\subsection*{Studi Kasus 2: Verifikasi Spesifikasi Program (Precondition dan Postcondition)}

Spesifikasi program sering dinyatakan sebagai implikasi: \emph{jika} input memenuhi precondition $P$, \emph{maka} output memenuhi postcondition $Q$. Membuktikan kebenaran spesifikasi ini sama dengan membuktikan implikasi $P \to Q$; dalam praktik kita dapat menggunakan bukti langsung (asumsikan $P$, tunjukkan $Q$) atau kontraposisi \cite{stanford_cs103,berkeley_logic}. Untuk fungsi konkret, kita dapat menulis assertion: untuk setiap input yang memenuhi $P$, periksa bahwa output memenuhi $Q$. Ini merupakan verifikasi berbasis tes yang mengilustrasikan ide bukti langsung.

\begin{contoh}
\textbf{Verifikasi precondition--postcondition dengan assertion}

Fungsi \code{abs\_kuadrat} mengembalikan $|n|^2$ untuk bilangan bulat $n$. Spesifikasi: jika $n$ bilangan bulat (precondition), maka hasil non-negatif dan $(\text{hasil}) = n^2$ (postcondition). Kita verifikasi dengan assertion untuk beberapa nilai yang memenuhi precondition.

\begin{lstlisting}[style=pythonstyle, caption={Verifikasi spesifikasi: jika precondition maka postcondition}]
def abs_kuadrat(n):
    """Return n^2 (precondition: n integer)."""
    return n * n

# Spesifikasi: P = "n integer", Q = "result >= 0 and result == n^2"
# Bukti langsung: asumsikan P (n integer), tunjukkan Q
def precondition(n):
    return isinstance(n, int)

def postcondition(n, result):
    return result >= 0 and result == n * n

# Verifikasi untuk beberapa input yang memenuhi P
for n in [-3, -1, 0, 1, 5]:
    assert precondition(n), "Precondition: n harus integer"
    result = abs_kuadrat(n)
    assert postcondition(n, result), "Postcondition: result >= 0 dan result == n^2"
print("Verifikasi berhasil: untuk setiap n integer, postcondition terpenuhi.")
\end{lstlisting}
\end{contoh}

Kedua studi kasus di atas menunjukkan penerapan aturan inferensi (modus ponens, modus tollens) dalam sistem aturan serta penerapan ide bukti langsung dalam verifikasi spesifikasi program \cite{rosen2019,stanford_cs103,mit_6042}.
